%% ============================================================
%% preamble.tex — ARES 논문
%% 컴파일: xelatex main.tex  (또는 latexmk -xelatex main.tex)
%% 한글: kotex + fontspec, macOS 기본 탑재 폰트(Apple SD Gothic Neo)만 사용하여
%%       별도 폰트 설치 없이 컴파일되도록 구성.
%% 코드: listings (Pygments/shell-escape 불필요). 코드 블록은 ASCII만 사용.
%% ============================================================

%% ----- 한국어 -----
\usepackage{kotex}
\usepackage{fontspec}

%% macOS 기본 폰트로 한글+라틴 모두 처리 (포터블).
%% 다른 환경이라면 아래 폰트명만 바꾸면 된다.
%%   - 서체 변경 예: \setmainhangulfont{KoPubWorld Batang}
\setmainfont{Apple SD Gothic Neo}
\setsansfont{Apple SD Gothic Neo}
\setmonofont{Menlo}
\setmainhangulfont{Apple SD Gothic Neo}
\setsanshangulfont{Apple SD Gothic Neo}
\setmonohangulfont{Apple SD Gothic Neo}

%% ----- 레이아웃 -----
\usepackage[a4paper, top=25mm, bottom=25mm, left=27mm, right=27mm]{geometry}
\linespread{1.18}
\setlength{\parskip}{0.4em}
\setlength{\parindent}{1em}

%% ----- 수식 -----
\usepackage{amsmath, amssymb}

%% ----- 그래픽/다이어그램 -----
\usepackage{graphicx}
\usepackage{tikz}
\usetikzlibrary{arrows.meta, positioning, shapes.geometric, calc, fit, backgrounds}

%% ----- 표 -----
\usepackage{booktabs}
\usepackage{tabularx}
\usepackage{multirow}
\usepackage{array}
\usepackage{makecell}
\renewcommand\theadfont{\bfseries\sffamily\small}

%% ----- 색상 -----
\usepackage{xcolor}
\definecolor{deepblue}{RGB}{22, 56, 100}
\definecolor{skyblue}{RGB}{70, 130, 200}
\definecolor{accentcopper}{RGB}{190, 100, 45}
\definecolor{rulegray}{RGB}{205, 210, 218}
\definecolor{codebg}{RGB}{248, 249, 251}
\definecolor{codecomment}{RGB}{110, 120, 135}
\definecolor{codekw}{RGB}{40, 90, 160}
\definecolor{codestr}{RGB}{150, 80, 40}

%% ----- 코드 (listings, ASCII 전용) -----
\usepackage{listings}
%% listings에는 JavaScript 내장 정의가 없으므로 직접 정의한다.
\lstdefinelanguage{JavaScript}{
  keywords={await, async, const, let, var, function, return, if, else, for, while, new, of, in, typeof, this},
  keywordstyle=\color{codekw}\bfseries,
  ndkeywords={true, false, null, undefined, Math, encoder},
  ndkeywordstyle=\color{codekw},
  sensitive=true,
  comment=[l]{//},
  morecomment=[s]{/*}{*/},
  morestring=[b]',
  morestring=[b]",
}
\lstdefinestyle{aresstyle}{
  backgroundcolor=\color{codebg},
  basicstyle=\ttfamily\footnotesize,
  commentstyle=\color{codecomment}\itshape,
  keywordstyle=\color{codekw}\bfseries,
  stringstyle=\color{codestr},
  numberstyle=\tiny\color{codecomment},
  breaklines=true,
  breakatwhitespace=false,
  captionpos=b,
  frame=leftline,
  framerule=1pt,
  rulecolor=\color{accentcopper},
  xleftmargin=10pt,
  framexleftmargin=10pt,
  showstringspaces=false,
  tabsize=2,
  numbers=none,
  upquote=true,
  columns=fullflexible,
}
\lstset{style=aresstyle}

%% ----- 캡션 -----
\usepackage{caption}
\captionsetup{font={small}, labelfont={small,bf,color=deepblue}}
\captionsetup[table]{position=above, skip=4pt}

%% ----- 섹션 제목 스타일 -----
\usepackage{titlesec}
\titleformat{\section}
  {\sffamily\Large\bfseries\color{deepblue}}
  {\thesection}{0.6em}{}
  [\vspace{2pt}{\color{rulegray}\hrule height 0.6pt}]
\titleformat{\subsection}
  {\sffamily\large\bfseries\color{accentcopper}}
  {\thesubsection}{0.6em}{}
\titleformat{\subsubsection}
  {\sffamily\normalsize\bfseries\color{deepblue}}
  {\thesubsubsection}{0.6em}{}
\titlespacing*{\section}{0pt}{16pt plus 2pt}{8pt}

%% ----- 유틸리티 -----
\usepackage{enumitem}
\setlist{itemsep=2pt, topsep=3pt, parsep=0pt}
\usepackage{float}
\usepackage{microtype}

%% ----- 한국어 캡션명 -----
\renewcommand{\figurename}{그림}
\renewcommand{\tablename}{표}
\renewcommand{\refname}{참고문헌}
\renewcommand{\contentsname}{목차}
\renewcommand{\abstractname}{초록}

%% ----- 하이퍼링크 -----
\usepackage[hidelinks]{hyperref}
\hypersetup{
  colorlinks=true,
  linkcolor=deepblue,
  citecolor=accentcopper,
  urlcolor=deepblue,
  pdftitle={ARES: Web Bluetooth 기반 무설치 교육용 자율 로버 플랫폼},
  pdfkeywords={ARES, Web Bluetooth, BLE, Blockly, 라즈베리파이 피코, 피지컬 컴퓨팅, 블록코딩}
}

%% 인용 마커 (수동 thebibliography 사용 시 \cite 동작)
%% 본문 강조용 헬퍼
\newcommand{\code}[1]{\texttt{\small #1}}
\newcommand{\keyterm}[1]{\textbf{#1}}
